\section{Nonlinear pairwise comparison on \texorpdfstring{$\SU\simeq S^3$}{SU(2) = S3}}
\label{sec:nonlinear}

Assume a Fr\'echet minimizer $C_*$ has been selected and
\[
  d(C_*,S_i)\le r<\frac\pi2
\]
for every $i$.  Put $v_i=\log_{C_*}(S_i)$.  Since no $S_i$ lies in the cut
locus and $C_*$ is stationary,
\[
  \sum_{i=1}^m v_i=0,
  \qquad
  D=\sum_{i=1}^m\norm{v_i}^2.
\]
Define the normalized pairwise dispersion
\[
  P_{\mathrm{pair}}=\frac1m\sum_{1\le i<j\le m}d(S_i,S_j)^2,
  \qquad s_r=\frac{\sin r}{r}.
\]

\begin{theorem}[Spherical dispersion comparison]
\label{thm:sphere}
Under the preceding hypotheses,
\[
  s_r^2D\le P_{\mathrm{pair}}\le D.
\]
The constants tend to equality as $r\downarrow0$.
\end{theorem}

The upper bound expresses the contraction of distance under the exponential
map on a positively curved sphere.  For the lower bound, the ball
$B(C_*,r)$ is strongly geodesically convex for $r<\pi/2$, and the differential
of the logarithm has norm at most $1/s_r$ there.  Combining these comparisons
with the Euclidean variance identity
\[
 \frac1m\sum_{i<j}\norm{v_i-v_j}^2
 =\sum_i\norm{v_i}^2
\]
proves the claim.  Full details are in Appendix~\ref{app:sphere}.

The theorem separates two effects.  Centering is exact: it is encoded by the
stationarity equation.  Curvature creates the only loss, through $s_r^2$.
For small $r$,
\[
 s_r^2=1-\frac{r^2}{3}+O(r^4),
\]
so the nonlinear and tangent dispersions agree to relative order $r^2$.

We will also use a simple deterministic transfer statement.  Let
$x=(x_1,\ldots,x_m)$ be a linear surrogate in a common Euclidean color space,
$\bar x=m^{-1}\sum_i x_i$, and
\[
 Q(x)=\sum_i\norm{x_i-\bar x}^2.
\]
If its centered coordinate error obeys
\[
 \left(\sum_i\norm{x_i-\bar x-v_i}^2\right)^{1/2}\le\sqrt m\,\varepsilon,
\]
then the reverse triangle inequality gives
\begin{equation}
\label{eq:det-transfer}
  \bigl|\sqrt{Q(x)}-\sqrt D\bigr|\le\sqrt m\,\varepsilon.
\end{equation}
Thus $D\ge m\eta^2$ implies
$Q\ge m(\eta-\varepsilon)_+^2$.  Equation~\eqref{eq:det-transfer} is exact
once its remainder premise is proved; deriving such a premise uniformly under
an interacting nonlinear gauge measure is outside the present theorem.
